\section{Introduction}
\label{sec:introduction}

\subsection{Motivation}
In modern industrial cyber-physical infrastructure---ranging from civil aviation propulsion systems to renewable power generation fleets---unexpected component failures incur catastrophic economic losses, unplanned operational downtime, and severe safety hazards. The proliferation of Industrial Internet of Things (IIoT) architectures has facilitated continuous, high-frequency telemetry acquisition across fleets of rotating machinery. Consequently, Predictive Maintenance (PdM) and Prognostics and Health Management (PHM) have shifted from retrospective breakdown repairs and calendar-based preventative servicing toward continuous condition-based prognostics \cite{saxena2008damage, ramasso2014performance}.

The core operational objective of predictive maintenance is twofold:
\begin{enumerate}
    \item \textbf{Early Anomaly Warning}: Detecting the earliest physical transition from nominal, healthy operation to an anomalous, degrading state, providing maintenance operators with sufficient warning lead time while constraining false alarm rates.
    \item \textbf{Prognostic Horizon Forecasting}: Accurately estimating the Remaining Useful Life (RUL) of the asset, defined as the number of operational cycles or operating hours remaining before component failure.
\end{enumerate}

The canonical benchmark for evaluating data-driven PdM methodologies is the NASA Commercial Modular Aero-Propulsion System Simulation (C-MAPSS) dataset suite \cite{saxena2008damage}. C-MAPSS provides high-dimensional run-to-failure telemetry across four subsets (FD001 through FD004), encompassing single and multi-regime operating conditions as well as single and dual fault modes.

\subsection{Problem Statement and Evaluation Vulnerabilities}
Despite a decade of literature applying deep neural networks, ensemble tree methods, and recurrent architectures to C-MAPSS, published evaluations suffer from pervasive methodological vulnerabilities that threaten real-world validity:
\begin{enumerate}
    \item \textbf{Circular Warning Lead Times}: In supervised classification studies, anomaly labels are commonly assigned by designating cycles where $RUL \le N$ (typically $N=30$) as ``anomalous''. When an evaluation claims an early warning lead time of 28 cycles before failure, this metric is fundamentally circular: it reflects the arbitrary choice of $N=30$ rather than the physical onset of degradation.
    \item \textbf{Optimistic Baseline Evaluation}: Many comparative studies benchmark complex deep architectures against uncalibrated or improperly tuned statistical baselines, failing to condition control charts on operating regimes or evaluating baselines under mismatched false alarm rates.
    \item \textbf{Cross-Condition Telemetry Leakage}: In multi-condition datasets (FD002, FD004), operating regimes dramatically distort sensor telemetry. In literature, normalization statistics are frequently fit across the entire dataset or fit on the target domain during cross-condition transfer, masking catastrophic domain shift.
    \item \textbf{Lack of Rigorous Statistical Correction}: Multiple comparative hypothesis tests are frequently published without Family-Wise Error Rate (FWER) corrections, leading to spurious claims of algorithmic superiority.
\end{enumerate}

\subsection{Research Questions}
To address these fundamental challenges, this paper investigates four central research questions:
\begin{description}
    \item[\textbf{RQ1}] \textbf{Prognostic and Diagnostic Accuracy}: How accurately can machine learning architectures identify anomalies and forecast degradation (RUL) across single-condition and multi-condition industrial telemetry fleets?
    \item[\textbf{RQ2}] \textbf{Machine Learning vs. Statistical Baselines}: How do machine learning models perform relative to condition-aware moving-average (EWMA) and static-threshold ($k$-sigma) baselines when evaluated at strictly equal false alarm rates across full trajectories?
    \item[\textbf{RQ3}] \textbf{Achievable Early Warning Lead Time}: What maximum warning lead time is reliably achievable in industrial telemetry fleets, at what false alarm rate, and how does supervised label definition ($N$) constrain achievable lead times?
    \item[\textbf{RQ4}] \textbf{Optimal Preprocessing and Dimensionality Reduction}: Which signal preprocessing techniques (temporal smoothing) and feature representations (raw flattened windows vs. PCA dimensionality reduction vs. summary statistics engineering) yield superior prognostic and diagnostic performance?
\end{description}

\subsection{Contributions}
The key contributions of this work are as follows:
\begin{itemize}
    \item \textbf{Leakage-Free Standardized Benchmark Pipeline}: A fully reproducible evaluation architecture implementing 5-fold shuffled \texttt{GroupKFold} by engine unit across 5 random seeds (0--4), guaranteeing zero test engine data leakage and seed-varying fold splits.
    \item \textbf{Unconstrained Lead-Time Operating Curves}: Sweep-based evaluation of warning lead time vs. false alarms per 100 healthy cycles across 7 methods (Static Threshold, EWMA, Mahalanobis Health Index, Isolation Forest, Random Forest, XGBoost, and LSTM), establishing reliable lead times at controlled industrial false alarm budgets ($\text{FA} \in \{0.1, 0.5, 1.0\}$).
    \item \textbf{Empirical Proof of Supervised Circularity}: A systematic label sensitivity sweep across $N \in \{20, 30, 50, 75, 100\}$ demonstrating that supervised alarm lead time tracks $N$ rather than physical damage, while proving that expanding $N \ge 75$ causes false alarm rates to escalate to $25.7$ per 100 cycles.
    \item \textbf{Identification of the Baseline Warning Lead-Time Paradox}: Discovery that while ML models vastly outperform statistical baselines in window-level classification ($F_1 = 0.753$--$0.858$ vs. $0.336$--$0.671$), a tuned regime-normalized static threshold warns earlier than ML models on multi-fault/multi-regime fleets under minimal false alarms ($p = 0.0264$ on FD003).
    \item \textbf{Cross-Condition Generalization Quantification}: Rigorous evaluation of cross-condition transfer (FD001 $\to$ FD002/FD004, FD003 $\to$ FD004) with normalizers fit strictly on the training domain, revealing a $103\%$ increase in RUL RMSE and a collapse of $F_1$ from $0.797$ down to $0.099$.
    \item \textbf{Distribution-Free Uncertainty Quantification}: Implementation of split conformal prediction intervals calibrated on whole held-out train engines at $90\%$ target coverage, reporting empirical test coverage and interval widths.
    \item \textbf{Physical Explainability}: TreeSHAP feature attributions on XGBoost demonstrating $60\%$--$80\%$ consensus with the sensors driving statistical control charts, identifying High-Pressure Compressor (HPC) outlet static pressure (\texttt{s11}) as the primary wear indicator.
\end{itemize}

\subsection{Paper Organization}
The remainder of this paper is organized as follows: Section~\ref{sec:related_work} reviews related literature in prognostic modeling and baseline benchmarking. Section~\ref{sec:dataset} presents the NASA C-MAPSS dataset suite and condition-aware regime normalization. Section~\ref{sec:methodology} formalizes the predictive architectures, evaluation protocols, and conformal calibration. Section~\ref{sec:experiments} presents comprehensive empirical findings across all experimental phases. Section~\ref{sec:discussion} analyzes the baseline paradox, lead-time circularity, and operational implications. Section~\ref{sec:limitations} evaluates methodological limitations, and Section~\ref{sec:conclusion} concludes.
